\documentclass[fleqn]{article} 
\usepackage[a3paper,landscape,vmargin=1.0cm,hmargin=1.5cm]{geometry} 

\usepackage{latexsym}
\usepackage{graphicx}
\usepackage{latexsym}

\usepackage{multicol}

\usepackage{epsfig}
\usepackage[dvips,xdvi]{color}

\usepackage{amsmath}
\usepackage{amssymb}

%--------------------------------------------------------------
	\setlength{\textwidth}{1120pt}
	\setlength{\textheight}{790pt}
%--------------------------------------------------------------
	\setlength{\oddsidemargin}{-50pt}
	\setlength{\evensidemargin}{-50pt}
%--------------------------------------------------------------
	\setlength{\topmargin}{-70pt}
	\setlength{\headheight}{0pt}
	\setlength{\topskip}{0pt}
%--------------------------------------------------------------
	\setlength{\columnseprule}{0.5pt}
	\setlength{\columnsep}{10pt}
%--------------------------------------------------------------
	\renewcommand{\baselinestretch}{1}
%--------------------------------------------------------------
	\setlength{\mathindent}{0.5cm}
	\setlength{\parindent}{7pt}
%--------------------------------------------------------------
	\renewcommand{\bottomfraction}{0.99}
	\renewcommand{\topfraction}{0.99}
	\renewcommand{\textfraction}{0.01}
%--------------------------------------------------------------
	\pagestyle{empty}
%--------------------------------------------------------------


%----------------------------------------------------------------
%	definition
%----------------------------------------------------------------
%----------------------------------------------------------------
%   equation 環境
%----------------------------------------------------------------
%		\def\topmat{\ \\ \vspace{-4mm}}
%		\def\botmat{\ \\ \vspace{-2mm}}
%		\def\bes{\vspace*{2mm}}
%		\def\ees{\ \\[-2mm]}
	%........................................................
	% equationの環境設定(\be,\ee)
	%      例 \bea 式 & 式 & 式 \nonumber\\ 式 & 式 & 式 \eea
	%........................................................
%		\def\be{\\ \ \\[-12mm] \begin{equation}}
%		\def\ee{\end{equation} \ \\[-11mm]}
		\newcommand{\be}{\begin{equation}}
		\newcommand{\ee}{\end{equation}}
	%........................................................
	% eqnarrayの環境設定(\bea,\eea)
	%      例 \bea 式 & 式 & 式 \nonumber\\ 式 & 式 & 式 \eea
	%........................................................
%		\def\bea{\\ \ \\[-16mm] \begin{eqnarray}}
%		\def\eea{\end{eqnarray} \ \\[-13mm]}
		\newcommand{\bea}{\begin{eqnarray}}
		\newcommand{\eea}{\end{eqnarray}}
	%........................................................
	% 番号なしeqnの環境設定(\nbe,\nee)
	%      例 \nbe 式 \\ 式  \nee
	%........................................................
%		\def\nbe{\\ \ \\[-12mm] \[}
%		\def\nee{\] \ \\[-11mm]}
		\newcommand{\nbe}{\[}
		\newcommand{\nee}{\]}
	%........................................................
	% matrixの入力(\bmatrix,\ematrix)
	%	例 \bmatrix{cc} 1&2\\3&4 \ematrix
	%........................................................
		\renewcommand{\bmatrix}{\left[\begin{array}}
		\newcommand{\ematrix}{\end{array}\right]}
	%........................................................
	% Greek character
	%........................................................
		\def\gggm{\mu}
		\def\ggge{\varepsilon}
		\def\gggw{\omega}
		\def\gggh{\theta}
		\def\ggga{\alpha}
		\def\gggb{\beta}
		\def\gggc{\gamma}
		\def\gggd{\delta}
		\def\gggp{\phi}
	%........................................................
	% others
	%........................................................
%		\def\half{\frac{1}{2}}
		\def\defeq{\stackrel{\rm def}{=}}


\renewcommand{\baselinestretch}{0.8}



%================================================================
\begin{document}
%================================================================
\rightline{
線形制御理論（2024.12.03)
\hspace{3cm}
\underline{学籍番号：　　　　　　　　　　　}
\hspace{1cm}
\underline{氏名：　　　　　　　　　　　　　}
\hspace{2cm}
}
%================================================================
\begin{multicols}{4}
%================================================================
\begin{enumerate}
%================================================================





%================================================================
\item
次の状態方程式で表される制御対象に対して一定値目標値に定常偏差なく追従するサーボ系をつくりたい．
\begin{align*}
	\frac{dx}{dt} & = 2 \cdot x + 1 \cdot u \\
	y &= 1 \cdot x
\end{align*}
以下の問いに答えよ．
\begin{enumerate}
\item
サーボ系を設計するために必要となる積分器の状態方程式を記せ．
\vspace{4cm}

\item
閉ループ系を安定化するコントローラを設計するために，制御対象と積分器をまとめて一つの状態方程式で表せ（拡大系）．
\vspace{6cm}

\item
すべての閉ループ系の極が$-3$となるコントローラを設計せよ．
\vspace{9cm}

\end{enumerate}

%================================================================
\vfill\null
\columnbreak
%================================================================



%================================================================









%================================================================
\item
次の１次元システム
\vspace{-1mm}
\begin{align}
	\frac{dx}{dt} &= 2 \cdot x + 1 \cdot u
	\ \ \ \  (A=2, B=1) 
\nonumber
\vspace{-4mm}
\end{align}
に対して評価関数
%\vspace{-1mm}
\begin{align}
	J &= \int_0^\infty  x^2 + r \cdot u^2 dt
	\ \ \ \  (Q=1, R=r)
\nonumber
%\vspace*{-8mm}
\end{align}
を最小化する最適フィードバック
%\vspace{-1mm}
\begin{align}
	u=Fx
\nonumber
%\vspace{-8mm}
\end{align}
を求め，その性質について調べたい．
%\vspace{-2mm}

\begin{enumerate}

\item
最適制御を設計するためにリカッチ方程式の正定解を求めよ．
\vspace{12cm}


\item
最適フィードバック$u=Fx$を求めよ．
\vspace{4cm}

%================================================================
\vfill\null
\columnbreak
%================================================================

\item
最適フィードバック$u=Fx$を施した後の閉ループ系の状態方程式を求めよ．
\vspace{4cm}


\item
$r \rightarrow 0$としたときの「入力」と「閉ループ系の極」を求めよ．
\vspace{4cm}

\item
$r \rightarrow \infty$としたときの「入力」と「閉ループ系の極」を求めよ．
\vspace{4cm}

\item
$r \rightarrow 0$なる評価関数を用いるということはどのような制御系を設計することに対応するのか，簡潔に述べよ．
\vspace{5cm}

\item
$r \rightarrow \infty$なる評価関数を用いるということはどのような制御系を設計することに対応するのか，簡潔に述べよ．
\vspace{5cm}

\end{enumerate}
%================================================================

%================================================================
\vfill\null
\columnbreak
%================================================================




%================================================================
\item
次の微分方程式で表されるシステムを考える．
\[
	 \frac{d^3 y}{dt^3} + a_2 \frac{d^2 y}{dt^2} + a_1 \frac{d y}{dt} + a_0 \cdot y = b_0 \cdot u
\]
ここで$y$は出力，$u$は入力とし,$a_0,a_1,a_2,b_0$は定数とする．

このシステムの状態$x$を
\[
	x = \bmatrix{c} x_1 \\ x_2           \\ x_3             \ematrix
      = \bmatrix{c} y   \\ {dy}/{dt}     \\ {d^2 y}/{dt^2} \ematrix
%      = \bmatrix{c} y   \\ \frac{dy}{dt} \\ \frac{d^2 y}{dt^2} \ematrix
\]
として状態方程式を求めたい．
以下の問いに答えよ．
\begin{enumerate}
\item 
$x_1, x_2, x_3$をそれぞれ時間微分し，状態と入力の関数として表せ．
\vspace{10cm}

\item
システムの状態方程式を求めよ．
\vspace{9cm}



\end{enumerate}

%================================================================
%\vfill\null
%\columnbreak
%================================================================


%================================================================



%================================================================
\end{enumerate}
%================================================================
\end{multicols}
%================================================================
\end{document}
%================================================================
